Our approach tests whether trustworthiness dimensions are empirically independent by computing cross-benchmark correlations and evaluating whether failure patterns cluster into distinct modes. The core insight is statistical: if reliability, robustness, and fairness arise from independent mechanisms, pairwise correlations should be near-zero (null hypothesis); if shared root causes exist, correlations should exceed random chance with at least moderate effect sizes ($r > 0.3$). Observed correlations far stronger than this threshold ($r > 0.99$) would indicate either near-redundancy across benchmarks or insufficient measurement resolution with only 3 dimensions.

\subsection{Dataset and Benchmark Selection}

We collected public benchmark results for 20 LLMs spanning three model size strata: small ($<$1B parameters, $n=6$), medium (1-10B, $n=9$), and large ($>$10B, $n=5$). Models include diverse architectures (GPT series, LLaMA series, Claude series, Mistral, Phi, Gemma) to control for architecture-specific confounds. Stratification by parameter count enables testing whether correlation patterns persist across scales or reflect model size artifacts.

\textbf{Benchmarks:}
\begin{itemize}
\item \textbf{TrustfulQA} (reliability): 817 questions testing factual accuracy and resistance to common misconceptions. Score = \% correct responses.
\item \textbf{AdvBench} (robustness): Attack success rate measuring vulnerability to adversarial jailbreaks and perturbations. Score = \% of attacks the model successfully defends against.
\item \textbf{BOLD} (fairness): Demographic bias metrics measuring sentiment/regard disparities across groups. Score = parity-adjusted fairness score (higher = more fair).
\end{itemize}

All scores normalized to [0,1] range via min-max scaling for correlation analysis. We selected these three benchmarks due to (1) public results availability for 20+ models, (2) conceptual coverage of key trustworthiness dimensions, and (3) established validity in prior work. Broader coverage (5-10 benchmarks) was deferred to future work due to data gaps across model families.

\textbf{Design Rationale:} 3-benchmark evaluation imposes clustering constraints ($k \geq 3$ requires $n \geq 3$ samples) but provides sufficient statistical power for correlation analysis ($n=20$ models $\times$ 3 benchmarks = 60 observations across 3 correlation pairs). This tradeoff prioritizes testing correlation existence (prediction P1) over taxonomy validation (prediction P2), with clustering serving as exploratory analysis rather than confirmatory test.

\subsection{Correlation Analysis}

We computed pairwise Spearman rank correlations for all benchmark pairs across the 20-model corpus. Spearman correlation measures monotonic relationships without assuming linearity, making it robust to ceiling effects (where high-performing models cluster near 100\% accuracy) and non-normal distributions.

\textbf{Statistical Testing:}
\begin{enumerate}
\item \textbf{Permutation test} (1000 iterations): Generate null distribution by shuffling benchmark scores, compute correlations for each permutation, compare observed correlations to null distribution percentiles.
\item \textbf{Bonferroni correction}: Adjust significance threshold to $\alpha = 0.01 / 3 = 0.0033$ to control family-wise error rate across three pairwise comparisons.
\item \textbf{Effect size interpretation}: Following Cohen's conventions, $r > 0.1$ (small), $r > 0.3$ (medium), $r > 0.5$ (large). We hypothesized $r > 0.3$ for shared root causes; $r > 0.99$ would indicate near-redundancy.
\end{enumerate}

\textbf{Stratified Analysis:} To control for model size confounds, we computed correlations separately within small, medium, and large strata. If correlations remain significant within size-homogeneous groups, coupling cannot be attributed to size-related performance trends (larger models uniformly better across dimensions).

\subsection{Hierarchical Clustering}

We applied Ward linkage hierarchical clustering to the $3 \times 3$ correlation matrix (converted to distance matrix via $d = 1 - r$) to test whether failure patterns organize into distinct modes. Ward linkage minimizes within-cluster variance, appropriate for correlation-based distances where tight coupling (high $r$) should yield small distances.

\textbf{Cluster Quality Metrics:}
\begin{itemize}
\item \textbf{Silhouette score} (range -1 to 1): Measures separation quality. Score $> 0.5$ indicates well-separated clusters; 0.3-0.5 acceptable; $<0.3$ poorly separated. Computed for $k=2$ clusters ($k \geq 3$ not computable with $n=3$ benchmarks).
\item \textbf{Bootstrap resampling} (1000 iterations): Resample models with replacement, recompute clustering, measure consistency via co-occurrence matrix (\% of iterations where benchmark pairs cluster together).
\item \textbf{Cophenetic correlation} (target $>0.7$): Measures how well dendrogram preserves pairwise distances from original correlation matrix.
\end{itemize}

\textbf{Success Criteria:} Silhouette $> 0.5$ AND bootstrap consistency $\geq 80\%$ would validate stable, well-separated clusters interpretable as distinct failure modes. Silhouette $< 0.5$ despite high bootstrap consistency would indicate stable but poorly separated groups---clusters reproducible but not semantically distinct.

\textbf{Methodological Constraint:} With only 3 benchmarks, we can evaluate $k=2$ clusters (sklearn silhouette requires $k <$ n\_samples). Testing $k=3-5$ failure modes (original hypothesis) requires expanding to $\geq 5$ benchmarks, motivating future work.

\subsection{Scale Invariance Testing}

To test whether failure mode clusters persist across model scales (prediction P3), we planned:
\begin{enumerate}
\item Compute separate correlation matrices within each size stratum (small/medium/large).
\item Apply hierarchical clustering to each stratum-specific matrix.
\item Use Mantel test to compare correlation matrices across strata pairs (small-medium, medium-large, small-large).
\item Success criterion: Mantel $r > 0.7$ for all pairs, indicating structure preserved across scales.
\end{enumerate}

\textbf{Note:} This analysis was blocked by clustering quality failure (h-m1 MUST\_WORK gate). However, stratified correlation analysis (Section 3.2) provides partial evidence: correlations remain $r > 0.98$ within all strata, suggesting coupling generalizes across scales even without formal Mantel test.

\subsection{Why This Design Tests the Hypothesis}

The methodology directly operationalizes the core claim: if trustworthiness failures share root causes manifesting across dimensions, then:
\begin{itemize}
\item \textbf{Correlation existence} (P1): Benchmark scores should correlate beyond random chance ($r > 0.3$, $p < 0.01$).
\item \textbf{Cluster structure} (P2): Correlations should organize into stable, well-separated groups (silhouette $> 0.5$, bootstrap $\geq 80\%$).
\item \textbf{Scale invariance} (P3): Patterns should persist across model sizes (Mantel $r > 0.7$).
\end{itemize}

Conversely, if dimensions are independent:
\begin{itemize}
\item Correlations should be near-zero ($r \sim 0$, $p > 0.05$).
\item Clustering should fail or produce unstable groups (bootstrap $< 60\%$).
\item Patterns should vary across strata (Mantel $r < 0.5$).
\end{itemize}

Observed $r > 0.99$ (far exceeding $r > 0.3$ threshold) validates correlation existence but introduces interpretive tension: does near-perfect coupling reflect unified constructs (benchmarks measure same thing) or insufficient resolution (3 benchmarks cannot distinguish dimensions that 5-10 might)? Clustering analysis intended to resolve this---distinct clusters would support separate dimensions despite high correlation---but silhouette = 0.274 indicates our 3-benchmark design cannot achieve required separation quality.

\subsection{Limitations and Mitigation}

\textbf{L1: Sample Size (3 Benchmarks):} Clustering limited to $k=2$ evaluation; $k \geq 3$ requires $n \geq 3$ samples. Mitigated by treating clustering as exploratory rather than confirmatory, with primary inference based on correlation analysis ($n=20$ models provides adequate power).

\textbf{L2: High Correlation Penalty:} $r > 0.99$ produces minimal distance variation (0.003-0.007 range), yielding low silhouette scores even for stable clusters. Mitigated by reporting both silhouette (separation) and bootstrap consistency (stability) separately, showing they measure orthogonal properties.

\textbf{L3: Public Data Constraints:} Benchmark results aggregated from published leaderboards may reflect inconsistent evaluation protocols. Mitigated by using large model sample ($n=20$) where measurement noise averages out, and focusing on rank correlations (robust to scale differences).

These limitations inform future work (Section 7): expand to 5-10 benchmarks to test $k \geq 3$ clustering with adequate sample size, apply alternative metrics (Calinski-Harabasz, Davies-Bouldin) suited to high-correlation data, and conduct controlled evaluations with standardized protocols rather than aggregating public results.
